\section{Introduction}

Automated ML research pipelines fail when hypotheses violate feasibility constraints during phase transitions. When hypotheses pass early validation checks but reach Phase 4 implementation only to discover they require unavailable datasets or human raters, all prior work is lost. This failure mode --- constraint violations discovered late rather than caught early --- forces expensive pipeline restarts and wastes compute on infeasible research directions.

Problem: existing workflow systems rely on schema validation to check structured outputs at phase boundaries, but schema validation checks structure only. Gap: semantic constraints (keywords) and compositional rules (cross-field logic) cannot be expressed. Existing research workflow systems rely on schema validation to check structured outputs at phase boundaries. Schema validation (e.g., Pydantic typed fields) validates structure --- types are correct, required fields present --- but cannot express semantic or compositional constraints. Consider a research pipeline with feasibility constraints: ``no synthetic data generation,'' ``no human evaluation,'' ``standard datasets only,'' ``existing benchmarks only.'' Schema validation checks that \texttt{dataset\_name} is a string but cannot verify it doesn't contain keywords like ``synthetic'' or ``generated.'' It validates that both \texttt{dataset\_type} and \texttt{dataset\_name} fields exist but cannot enforce the compositional rule ``IF dataset\_type equals `standard' THEN dataset\_name must be in the approved standard datasets list.''

This expressiveness gap has consequences. Workflow orchestration frameworks like FlowXpert and LangChain focus on troubleshooting workflows after failures occur, not preventing failures proactively. Schema-only validation at phase boundaries misses semantic violations (keyword patterns) and compositional failures (cross-field logic, state-based checks), allowing infeasible hypotheses to propagate through expensive implementation phases before discovery. No existing research applies formal constraint validation methods to research pipeline phase transitions.

We address this gap with a novel insight: three-layer contract-based validation (schema + pattern + compositional contracts) forces explicit checking of constraints that schema-only validation cannot express. Each layer targets a distinct violation class: schemas validate structure (types, required fields), pattern validators check semantic content (forbidden keywords), and compositional contracts enforce cross-field logic (conditional dependencies) and state-based rules (membership in external registries). The layers are complementary, not redundant --- each catches violations the previous layer misses.

Building on this insight, we make the following contributions:

\begin{enumerate}
\item \textbf{First application of Design-by-Contract to semantic constraint validation in research workflows.} We adapt preconditions, postconditions, and invariants from code correctness (Eiffel, C++26 contracts) to constraint-preserving phase transitions in research pipelines, applying contract concepts to semantic and compositional constraints (not just task idempotence as in prior workflow research).

\item \textbf{Three-layer validation architecture.} We implement and validate a schema $\rightarrow$ pattern $\rightarrow$ contract hierarchy where each layer complements the previous, achieving 88\% violation detection at boundaries versus 40\% for schema-only (48 percentage point gap).

\item \textbf{Empirical validation on placeholder workflows.} We demonstrate 100\% downstream failure reduction on a 100-case placeholder hypothesis corpus with four feasibility constraints on placeholder content with typed interfaces (external validity to substantive research unproven, deferred to Phase 5), eliminating Phase 4/5 failures that occur when violations propagate undetected through phase boundaries.

\item \textbf{Contract specification templates for rapid deployment.} We provide reusable templates for semantic constraints (keyword blacklists), cross-field constraints (conditional membership), and state-based constraints (registry validation), enabling research teams to adopt contract-based validation with minimal implementation cost (25--27 lines of code for the contract layer).
\end{enumerate}

The remainder of this paper is organized as follows. Section 2 positions our work relative to workflow orchestration, formal methods, and multi-level validation research. Section 3 presents our three-layer validation methodology and design rationale. Section 4 describes our experimental setup with four sub-hypotheses testing existence, mechanism, and end-to-end effectiveness. Section 5 presents results demonstrating 100\% failure reduction. Section 6 discusses findings, limitations, and implications. Section 7 concludes with future directions for auto-generating contracts and extending validation to production ML pipelines.
